\documentclass[10pt,a4paper]{amsart}
\usepackage{amsmath,amssymb,amsthm,graphicx}
% Retained minimal preamble for the Stacks Project chapter sites-cohomology.tex (Cohomology on Sites).
% Verbatim excerpt of the chapter preamble: Xy-pic setup (native preamble lines 14-18) and the native
% theorem declarations with the mathematical macros (native preamble lines 154-217). Copied mechanically;
% nothing added, removed or redefined. The wrapper class and the amsmath/amssymb/amsthm/graphicx base
% packages are supplied by the assembler, so the class/externaldocument part of the native preamble is omitted.
% For commutative diagrams we use Xy-pic
\usepackage[all]{xy}

% We use 2cell for 2-commutative diagrams.
\xyoption{2cell}
% Theorem environments.
%
\theoremstyle{plain}
\newtheorem{theorem}[subsection]{Theorem}
\newtheorem{proposition}[subsection]{Proposition}
\newtheorem{lemma}[subsection]{Lemma}

\theoremstyle{definition}
\newtheorem{definition}[subsection]{Definition}
\newtheorem{example}[subsection]{Example}
\newtheorem{exercise}[subsection]{Exercise}
\newtheorem{situation}[subsection]{Situation}

\theoremstyle{remark}
\newtheorem{remark}[subsection]{Remark}
\newtheorem{remarks}[subsection]{Remarks}

\numberwithin{equation}{subsection}

% Macros
%
\def\lim{\mathop{\mathrm{lim}}\nolimits}
\def\colim{\mathop{\mathrm{colim}}\nolimits}
\def\Spec{\mathop{\mathrm{Spec}}}
\def\Hom{\mathop{\mathrm{Hom}}\nolimits}
\def\Ext{\mathop{\mathrm{Ext}}\nolimits}
\def\SheafHom{\mathop{\mathcal{H}\!\mathit{om}}\nolimits}
\def\SheafExt{\mathop{\mathcal{E}\!\mathit{xt}}\nolimits}
\def\Sch{\mathit{Sch}}
\def\Mor{\mathop{\mathrm{Mor}}\nolimits}
\def\Ob{\mathop{\mathrm{Ob}}\nolimits}
\def\Sh{\mathop{\mathit{Sh}}\nolimits}
\def\NL{\mathop{N\!L}\nolimits}
\def\CH{\mathop{\mathrm{CH}}\nolimits}
\def\proetale{{pro\text{-}\acute{e}tale}}
\def\etale{{\acute{e}tale}}
\def\QCoh{\mathit{QCoh}}
\def\Ker{\mathop{\mathrm{Ker}}}
\def\Im{\mathop{\mathrm{Im}}}
\def\Coker{\mathop{\mathrm{Coker}}}
\def\Coim{\mathop{\mathrm{Coim}}}

% Boxtimes
%
\DeclareMathSymbol{\boxtimes}{\mathbin}{AMSa}{"02}

%
% Macros for moduli stacks/spaces
%
\def\QCohstack{\mathcal{QC}\!\mathit{oh}}
\def\Cohstack{\mathcal{C}\!\mathit{oh}}
\def\Spacesstack{\mathcal{S}\!\mathit{paces}}
\def\Quotfunctor{\mathrm{Quot}}
\def\Hilbfunctor{\mathrm{Hilb}}
\def\Curvesstack{\mathcal{C}\!\mathit{urves}}
\def\Polarizedstack{\mathcal{P}\!\mathit{olarized}}
\def\Complexesstack{\mathcal{C}\!\mathit{omplexes}}
% \Pic is the operator that assigns to X its picard group, usage \Pic(X)
% \Picardstack_{X/B} denotes the Picard stack of X over B
% \Picardfunctor_{X/B} denotes the Picard functor of X over B
\def\Pic{\mathop{\mathrm{Pic}}\nolimits}
\def\Picardstack{\mathcal{P}\!\mathit{ic}}
\def\Picardfunctor{\mathrm{Pic}}
\def\Deformationcategory{\mathcal{D}\!\mathit{ef}}
\usepackage{cleveref}
\crefname{originalreference}{original reference}{original references}
\Crefname{originalreference}{Original reference}{Original references}
\makeatletter
\newlabel{lemma-tor-amplitude-pullback}{{08H5}{source}{}{sourcepointer-1}{}}
\newlabel{lemma-tor-amplitude-pullback@cref}{{[originalreference][0][]08H5}{[1][0][]source}{}{}{}}
\makeatother
\crefname{originallemma}{original lemma}{original lemmas}
\Crefname{originallemma}{Original lemma}{Original lemmas}
\makeatletter
\newlabel{sites-modules-lemma-check-exactness-stalks}{{05V3}{source}{}{sourcepointer-2}{}}
\newlabel{sites-modules-lemma-check-exactness-stalks@cref}{{[originallemma][0][]05V3}{[1][0][]source}{}{}{}}
\makeatother
\title{Tor amplitude can be checked on the stalks when the site has enough points, and always passes to stalks}
\author{The Stacks Project Authors}
\date{stacks-project:sites-cohomology.tex:local-snapshot-20261011; GNU FDL 1.2 or later}
\begin{document}
\maketitle
\noindent\textit{Editorial note.} All delivered mathematical text is the author's own, copied from the pinned Stacks Project chapter file sites-cohomology.tex; no statement or proof step was rewritten, reordered or removed.\par\medskip
\noindent\textit{Editorial note.} The counted proof is the author's complete argument: the proof is written out in full for both parts: the pullback interpretation of the stalk and the stalk criterion for the converse.\par\medskip
\noindent\textit{Editorial note.} The label lemma-tor-amplitude-stalk names the counted statement, and the proof environment immediately following it in the source is delivered verbatim as its proof.\par\medskip
\noindent\textit{Editorial note.} The counted statement is carried by the author's own lemma environment in the chapter source.\par\medskip
\noindent\textit{Editorial note.} The statement and proof are copied from the pinned chapter file; the two cited inputs are kept as external pointers with their canonical Stacks tags.\par\medskip
\noindent\textit{Allowed standard prerequisites (editorial).}\begin{itemize}
\item Tor amplitude and its behaviour under pullback. Same chapter, the definition of tor amplitude in an interval and the lemma stating that it is preserved by pullback along a morphism of ringed sites. Supplies the reduction in part (1), where taking the stalk at a point is a pullback of ringed sites.
\item Exactness can be checked on stalks. Modules on Sites, the lemma allowing the vanishing of the cohomology of a tensor product to be tested at the points of the site. Supplies the converse in part (2) when the site has enough points.
\end{itemize}\medskip
\ifdefined\hypertarget\hypertarget{sourcepointer-1}{}\fi
\noindent\textit{Source reference.} Original reference 08H5: sources/stacks-sites-cohomology/pinned-original-source.tex, source lines 12543--12543. Prerequisite pointer invoked at proof line(s) 12706; it names a standard fact of the same chapter that the counted argument uses, and every step of the counted proof built on it is retained verbatim in the delivered body.\par\medskip
\ifdefined\hypertarget\hypertarget{sourcepointer-2}{}\fi
\noindent\textit{Source reference.} Original lemma 05V3: sources/stacks-chow/tags.txt, source lines 7537--7537. Prerequisite pointer invoked at proof line(s) 12713; it names a standard fact of a sibling chapter that the counted argument uses, and every step of the counted proof built on it is retained verbatim in the delivered body.\par\medskip
\begin{lemma}
\label{lemma-tor-amplitude-stalk}
Let $(\mathcal{C}, \mathcal{O})$ be a ringed site.
Let $E$ be an object of $D(\mathcal{O})$.
Let $a, b \in \mathbf{Z}$.
\begin{enumerate}
\item If $E$ has tor amplitude in $[a, b]$, then for every point $p$
of the site $\mathcal{C}$ the object $E_p$ of $D(\mathcal{O}_p)$
has tor amplitude in $[a, b]$.
\item If $\mathcal{C}$ has enough points, then the converse is true.
\end{enumerate}
\end{lemma}

\begin{proof}
Proof of (1). This follows because taking stalks at $p$ is
the same as pulling back by the morphism of ringed sites
$(p, \mathcal{O}_p) \to (\mathcal{C}, \mathcal{O})$ and hence
we can apply Lemma \ref{lemma-tor-amplitude-pullback}.

\medskip\noindent
Proof of (2). If $\mathcal{C}$ has enough points, then we can check
vanishing of
$H^i(E \otimes_\mathcal{O}^\mathbf{L} \mathcal{F})$
at stalks, see
Modules on Sites, Lemma \ref{sites-modules-lemma-check-exactness-stalks}.
Since $H^i(E \otimes_\mathcal{O}^\mathbf{L} \mathcal{F})_p =
H^i(E_p \otimes_{\mathcal{O}_p}^\mathbf{L} \mathcal{F}_p)$ we conclude.
\end{proof}

\end{document}
